%! Unit 5: Radical Functions — Summative Assessment — Unit Test (KEY)
\documentclass[10pt]{article}
\usepackage{algebra2-article}
\usepackage{algebra2-key}   % pulls in -boxes; do NOT also load -boxes
\usepackage{grading}
\gradingassignment{alg2-u05-actual-test}
\gradingcourse{algebra_2}
\usepackage{algebra2-graded}
% Coordinate grid: {xmin}{xmax}{ymin}{ymax}
\newcommand{\lgrid}[4]{%
  \draw[step=1, linegray!40, very thin] (#1,#3) grid (#2,#4);
  \draw[->, semithick, charcoal] (#1,0)--(#2,0) node[below right, font=\scriptsize]{$x$};
  \draw[->, semithick, charcoal] (0,#3)--(0,#4) node[left, font=\scriptsize]{$y$};}
\newcommand{\dpt}[2]{\fill[charcoal] (#1,#2) circle (2.6pt);}

% Part-divider strip (headlinebox takes one arg: the background color)
\newcommand{\parthead}[1]{%
  \vspace{6pt}\noindent
  \begin{headlinebox}{forest}\color{white}\bfseries #1\end{headlinebox}%
  \vspace{4pt}}

\begin{document}

\pageheader{Unit 5: Radical Functions}{Unit Test --- Answer Key}

\vspace{0.06in}
\noindent\textbf{Instructions:} Show all work. No notes unless your teacher says so.

% ======================================================================
\parthead{Part A --- Vocabulary (8 pts)}
Write the letter of the best description next to each term.

\vspace{4pt}
\noindent
\begin{minipage}[t]{0.36\textwidth}
\begin{enumerate}[leftmargin=2.2em, itemsep=7pt, label=\arabic*.]
  \item Index \inblank{\blank{1.2cm}}{\gradedblank[id=q1, type=choice, points=1, answer={E}]{1.2cm-2pt}{12pt}}
  \item Radicand \inblank{\blank{1.2cm}}{\gradedblank[id=q2, type=choice, points=1, answer={G}]{1.2cm-2pt}{12pt}}
  \item Principal root \inblank{\blank{1.2cm}}{\gradedblank[id=q3, type=choice, points=1, answer={B}]{1.2cm-2pt}{12pt}}
  \item Conjugate \inblank{\blank{1.2cm}}{\gradedblank[id=q4, type=choice, points=1, answer={D}]{1.2cm-2pt}{12pt}}
  \item Extraneous solution \inblank{\blank{1.2cm}}{\gradedblank[id=q5, type=choice, points=1, answer={A}]{1.2cm-2pt}{12pt}}
  \item Composition \inblank{\blank{1.2cm}}{\gradedblank[id=q6, type=choice, points=1, answer={H}]{1.2cm-2pt}{12pt}}
  \item Inverse function \inblank{\blank{1.2cm}}{\gradedblank[id=q7, type=choice, points=1, answer={C}]{1.2cm-2pt}{12pt}}
  \item Anchor point \inblank{\blank{1.2cm}}{\gradedblank[id=q8, type=choice, points=1, answer={F}]{1.2cm-2pt}{12pt}}
\end{enumerate}
\end{minipage}%
\hfill
\begin{minipage}[t]{0.61\textwidth}
\small
\begin{description}[leftmargin=1.4em, itemsep=3pt]
  \item[(A)] A candidate that satisfies the powered-up equation but fails when it is tested in the original.
  \item[(B)] The non-negative value the radical symbol by itself returns; a $\pm$ must be written in by hand.
  \item[(C)] The function that reverses another's inputs and outputs; its graph is the mirror image over the line $y=x$.
  \item[(D)] A binomial with the same two terms but the opposite middle sign; multiplying by it turns an irrational denominator rational.
  \item[(E)] The small number in the crook of the radical symbol that says which power is being undone.
  \item[(F)] The point $(h,k)$ that positions a transformed radical graph --- an endpoint when the index is even, a turning point when it is odd.
  \item[(G)] Everything written underneath the radical symbol; when the index is even, it may not be negative.
  \item[(H)] The operation that feeds one function's output into another function's input, written $(f\circ g)(x)$.
\end{description}
\end{minipage}

% ======================================================================
\parthead{Part B --- Multiple Choice (12 pts)}
Circle the best answer.

\begin{enumerate}[leftmargin=1.9em, itemsep=9pt]
  \item Written with a rational exponent, $\sqrt[5]{x^4}$ is\rlap{\hspace{4pt}\choicebox{\gradedblank[id=q9, type=choice, points=2, answer={B}]{0.3in}{10pt}}}
    \hfill (A) $x^{5/4}$ \quad \textcolor{keyred}{\underline{(B) $x^{4/5}$}} \quad (C) $x^{20}$ \quad (D) $x^{-5/4}$
  \item The domain of $f(x)=\sqrt{9-3x}$ is\rlap{\hspace{4pt}\choicebox{\gradedblank[id=q10, type=choice, points=2, answer={A}]{0.3in}{10pt}}}
    \hfill \textcolor{keyred}{\underline{(A) $x\le 3$}} \quad (B) $x\ge 3$ \quad (C) $x\ge -3$ \quad (D) all real numbers
  \item The graph of $y=\sqrt{2x-10}+3$ begins at the point\rlap{\hspace{4pt}\choicebox{\gradedblank[id=q11, type=choice, points=2, answer={B}]{0.3in}{10pt}}}
    \hfill (A) $(10,3)$ \quad \textcolor{keyred}{\underline{(B) $(5,3)$}} \quad (C) $(5,-3)$ \quad (D) $(-5,3)$
  \item In simplest form, $\sqrt{98}+\sqrt{32}$ is\rlap{\hspace{4pt}\choicebox{\gradedblank[id=q12, type=choice, points=2, answer={D}]{0.3in}{10pt}}}
    \hfill (A) $\sqrt{130}$ \quad (B) $28\sqrt2$ \quad (C) $11\sqrt{130}$ \quad \textcolor{keyred}{\underline{(D) $11\sqrt2$}}
  \item \letterbox{\gradedblank[id=q13, type=choice, points=2, answer={B}]{0.4in}{10pt}}Which equation has \emph{no real solution}?
    \\*[3pt]\hspace*{1.2em} (A) $\sqrt[3]{x+2}=-3$
    \hspace{1.0cm} \textcolor{keyred}{\underline{(B) $\sqrt{x+1}=-5$}}
    \hspace{1.0cm} (C) $\sqrt{x-4}=6$
    \hspace{1.0cm} (D) $\sqrt[3]{x}=0$
  \item If $f(x)=\sqrt{x}$ and $g(x)=x+4$, then $(f\circ g)(x)=$\rlap{\hspace{3pt}\choicebox{\gradedblank[id=q14, type=choice, points=2, answer={A}]{0.26in}{10pt}}}
    \hfill \textcolor{keyred}{\underline{(A) $\sqrt{x+4}$}} \quad (B) $\sqrt{x}+4$ \quad (C) $\sqrt{4x}$ \quad (D) $x\sqrt{x}+4\sqrt{x}$
\end{enumerate}

\begin{teachernote}
Item 1: \textbf{denominator = index, numerator = power}. (A) is the signature Unit 5 error --- the
flipped fraction --- and (C) multiplies index by power. Item 2: set the \emph{radicand} $\ge0$:
$9-3x\ge0$ gives $-3x\ge-9$, and dividing by $-3$ \textbf{flips} the inequality. (B) is the
un-flipped answer. Item 3: factor first --- $\sqrt{2x-10}=\sqrt2\sqrt{x-5}$, so the graph starts at
$x=5$, not $x=10$; (A) is the students who never factored. Item 4: $7\sqrt2+4\sqrt2$; (A) adds
radicands, (B) multiplies the coefficients, (C) adds both. Item 5: an isolated \emph{even} root is
never negative. Distractor (A) is the one to sort for --- a student who learned ``roots can't be
negative'' instead of ``\emph{even} roots can't'' picks it, and that is a whole-unit failure, not a
slip; $\sqrt[3]{x+2}=-3$ solves fine ($x=-29$). Item 6: the function next to the $x$ runs first. (B)
is the order error, (D) reads $\circ$ as multiplication.
\end{teachernote}

% ======================================================================
\parthead{Part C --- Short Answer \& Computation (40 pts)}
Show your work. Assume every variable is non-negative unless told otherwise.

\begin{enumerate}[leftmargin=1.9em, itemsep=10pt]
  \item \textbf{Read the graph.} The radical function $f$ is graphed below; each gridline is one unit
        and the solid dot is the endpoint.
        \begin{center}
        \begin{tikzpicture}[scale=0.45]
          \lgrid{-6}{8}{-4}{4}
          \begin{scope}
            \clip (-6,-4) rectangle (8,4);
            \draw[very thick, forest, domain=-4:8, samples=120, smooth]
              plot (\x,{sqrt((\x)+4)-1});
          \end{scope}
          \fill[forest] (-4,-1) circle (4pt);
          \foreach \x in {-5,-2,2,4,6}
            \draw[charcoal] (\x,-0.15)--(\x,0.15) node[below=1pt, font=\tiny]{$\x$};
          \dpt{-3}{0}
          \dpt{0}{1}
          \dpt{5}{2}
        \end{tikzpicture}
        \end{center}
        (a) Coordinates of the endpoint: \inblank{\blank{2.4cm}}{\gradedblank[id=q15a, type=short, points=0.5, answer={(-4,-1)}, check=text]{2.4cm-2pt}{10pt}}
        \\[4pt]
        (b) Domain (interval notation): \inblank{\blank{3.0cm}}{\gradedblank[id=q15b, type=short, points=0.5, answer={[-4,inf)}, check=text, display={$[-4,\infty)$}]{3.0cm-2pt}{10pt}}; \quad range: \inblank{\blank{3.0cm}}{\gradedblank[id=q15c, type=short, points=0.5, answer={[-1,inf)}, check=text, display={$[-1,\infty)$}]{3.0cm-2pt}{10pt}}
        \\[4pt]
        (c) $x$-intercept: \inblank{\blank{2.0cm}}{\gradedblank[id=q15d, type=short, points=0.5, answer={(-3,0)}, check=text]{2.0cm-2pt}{10pt}}; \quad $y$-intercept: \inblank{\blank{2.0cm}}{\gradedblank[id=q15e, type=short, points=0.5, answer={(0,1)}, check=text]{2.0cm-2pt}{10pt}}
        \\[4pt]
        (d) Interval(s) where $f$ is increasing: \inblank{\blank{4.0cm}}{\gradedblank[id=q15f, type=short, points=0.5, answer={[-4,inf)}, check=text, display={$[-4,\infty)$ --- the whole graph}]{4.0cm-2pt}{10pt}}
        \\*[4pt]
        (e) Does $f$ have an absolute maximum or an absolute minimum? Give it and say where it occurs.
 \rubricbox{\gradedblank[id=q15g, type=rubric, points=1, rubric={../rubrics/alg2-u05-actual-test-q15g.yaml}, display={absolute \textbf{minimum} $-1$ at the endpoint $x=-4$; no maximum}]{\linewidth}{10pt}}\par
        (f) End behavior: as $x\to\infty$, $f(x)\to$ \inblank{\blank{2.0cm}}{\gradedblank[id=q15h, type=short, points=0.5, answer={inf}, check=text, display={$\infty$}]{2.0cm-2pt}{10pt}};
            \ as $x\to-\infty$, \rubricbox{\gradedblank[id=q15i, type=rubric, points=1, rubric={../rubrics/alg2-u05-actual-test-q15i.yaml}, display={nothing --- the graph \textbf{stops} at $x=-4$}]{\linewidth}{10pt}}\par
        (g) Interval(s) where $f(x)<0$: \inblank{\blank{4.0cm}}{\gradedblank[id=q15j, type=short, points=0.5, answer={[-4,-3)}, check=text]{4.0cm-2pt}{10pt}}

  \item \textbf{Radicals and rational exponents.} Rewrite or evaluate. Leave no negative exponents.
        \\[4pt]
        (a) Write $\sqrt[3]{x^2}$ with a rational exponent. \inblank{\blank{3.0cm}}{\gradedblank[id=q16a, type=short, points=1, answer={x^(2/3)}, display={$x^{2/3}$}]{3.0cm-2pt}{10pt}}
        \\[4pt]
        (b) Write $y^{7/2}$ as a radical. \inblank{\blank{3.0cm}}{\gradedblank[id=q16b, type=short, points=1, answer={sqrt(y^7)}, display={$\sqrt{y^7}$, or $\left(\sqrt{y}\right)^7$}]{3.0cm-2pt}{10pt}}
        \\[4pt]
        (c) $8^{2/3}=$ \inblank{\blank{2.2cm}}{\gradedblank[id=q16c, type=short, points=1, answer={4}, check=numeric]{2.2cm-2pt}{10pt}} \hspace{1.2cm}
        (d) $25^{-1/2}=$ \inblank{\blank{2.2cm}}{\gradedblank[id=q16d, type=short, points=1, answer={1/5}, check=numeric, display={$1/5$}]{2.2cm-2pt}{10pt}}
        \\[4pt]
        \answerbox[Answer:]{\gradedblank[id=q16e, type=short, points=1, answer={x^(5/4)}, display={$\sqrt[4]{x^5}$}]{1.1in}{10pt}}(e) Write $x^{3/4}\cdot x^{1/2}$ as a single radical.
        \\
        \ans{Add the exponents: $x^{3/4+1/2}=x^{3/4+2/4}=x^{5/4}=\sqrt[4]{x^5}$. As radicals the two
        factors could not be combined at all --- mismatched indices --- which is exactly why the
        rational-exponent notation exists.}

  \item \textbf{Simplify. Add or subtract where you can.}
        \\[6pt]
        \answerbox[(a)]{\gradedblank[id=q17a, type=short, points=1, answer={7sqrt(2)}, display={$7\sqrt2$}]{0.95in}{10pt}\quad\textbf{(b)}\ \gradedblank[id=q17b, type=short, points=1, answer={2*5^(1/3)}, display={$2\sqrt[3]{5}$}]{0.95in}{10pt}}(a) $\sqrt{98}$ \hspace{2.4cm} (b) $\sqrt[3]{40}$
        \\[4pt]
        \ans{(a) Largest perfect square: $\sqrt{49\cdot2}=7\sqrt2$.\\[3pt]
        (b) The index sets the group size, so hunt for a perfect \emph{cube}:
        $\sqrt[3]{8\cdot5}=2\sqrt[3]{5}$. (The $4$ inside is a perfect square, which buys nothing
        here.)}
        \\
        \answerbox[(c)]{\gradedblank[id=q17c, type=short, points=1, answer={3x^2sqrt(2x)}, display={$3x^2\sqrt{2x}$}]{0.95in}{10pt}\quad\textbf{(d)}\ \gradedblank[id=q17d, type=short, points=1, answer={sqrt(2)}, display={$\sqrt2$}]{0.95in}{10pt}}(c) $\sqrt{18x^5}$ \hspace{2.1cm} (d) $3\sqrt{8}-\sqrt{50}$
        \\[4pt]
        \ans{(c) $\sqrt{9x^4\cdot 2x}=3x^2\sqrt{2x}$ --- divide the exponent by the index: $x^5$ sends
        two $x$'s out and leaves one in.\\[3pt]
        (d) Simplify first, decide second: $3\cdot2\sqrt2-5\sqrt2=6\sqrt2-5\sqrt2=\sqrt2$. Before
        simplifying these look like unlike radicals.}

  \item \textbf{Multiply, then rationalize.} Every answer must be in simplest form.
        \\[6pt]
        \answerbox[(a)]{\gradedblank[id=q18a, type=short, points=1, answer={2sqrt(35)}, display={$2\sqrt{35}$}]{0.95in}{10pt}\quad\textbf{(b)}\ \gradedblank[id=q18b, type=short, points=1, answer={2}, check=numeric, display={$2$}]{0.95in}{10pt}}(a) $\sqrt{10}\cdot\sqrt{14}$ \hspace{1.8cm} (b) $\left(3+\sqrt7\right)\left(3-\sqrt7\right)$
        \\[4pt]
        \ans{(a) $\sqrt{140}=\sqrt{4\cdot35}=2\sqrt{35}$. Neither factor simplifies on its own --- the
        \emph{product} does, so the last step is the graded one.\\[3pt]
        (b) Conjugates: $9-7=2$, a rational number. Compare
        $\left(3+\sqrt7\right)^2=16+6\sqrt7$, where the radical survives.}
        \\
        \answerbox[(c)]{\gradedblank[id=q18c, type=short, points=1, answer={2sqrt(5)}, display={$2\sqrt5$}]{0.95in}{10pt}\quad\textbf{(d)}\ \gradedblank[id=q18d, type=short, points=2, answer={4+sqrt(10)}, display={$4+\sqrt{10}$}]{0.95in}{10pt}}(c) $\dfrac{10}{\sqrt5}$ \hspace{3.0cm} (d) $\dfrac{6}{4-\sqrt{10}}$
        \\[4pt]
        \ans{(c) Multiply by $\frac{\sqrt5}{\sqrt5}=1$: $\dfrac{10\sqrt5}{5}=2\sqrt5$.\\[3pt]
        (d) Multiply by the \emph{conjugate} $\frac{4+\sqrt{10}}{4+\sqrt{10}}$:
        $\dfrac{6\left(4+\sqrt{10}\right)}{16-10}=\dfrac{6\left(4+\sqrt{10}\right)}{6}=4+\sqrt{10}$.
        Using $4-\sqrt{10}$ again would only produce a new radical.}

  \boxguard[16]
  \item \textbf{Transformations.} Let $g(x)=-3\sqrt{x-2}+1$.
        \\*[4pt]
        (a) List every transformation that takes $y=\sqrt{x}$ to $g$.
 \rubricbox{\gradedblank[id=q19a, type=rubric, points=1, rubric={../rubrics/alg2-u05-actual-test-q19a.yaml}, display={right $2$, stretch $\times3$, reflect over $x$-axis, up $1$}]{\linewidth}{10pt}}\par
        (b) Anchor point: \inblank{\blank{2.2cm}}{\gradedblank[id=q19b, type=short, points=0.5, answer={(2,1)}, check=text]{2.2cm-2pt}{10pt}}
        \\[4pt]
        (c) Domain: \inblank{\blank{2.8cm}}{\gradedblank[id=q19c, type=short, points=0.5, answer={[2,inf)}, check=text, display={$[2,\infty)$}]{2.8cm-2pt}{10pt}}; \quad range: \inblank{\blank{2.8cm}}{\gradedblank[id=q19d, type=short, points=0.5, answer={(-inf,1]}, check=text, display={$(-\infty,1]$}]{2.8cm-2pt}{10pt}}
        \\[4pt]
        (d) Is the anchor an absolute maximum or an absolute minimum? \inblank{\blank{4.0cm}}{\gradedblank[id=q19e, type=short, points=0.5, answer={maximum}, check=text, display={maximum ($a<0$ turns the arm over)}]{4.0cm-2pt}{10pt}}
        \\[4pt]
        (e) A different graph of the form $y=\sqrt{x-h}+k$ begins at $(-4,3)$. Write its equation.
            \inblank{\blank{3.6cm}}{\gradedblank[id=q19f, type=short, points=1, answer={y=sqrt(x+4)+3}, display={$y=\sqrt{x+4}+3$}]{3.6cm-2pt}{10pt}}
        \\*[4pt]
        (f) For $p(x)=\sqrt[3]{x+5}-1$, give the anchor point \inblank{\blank{2.2cm}}{\gradedblank[id=q19g, type=short, points=0.5, answer={(-5,-1)}, check=text]{2.2cm-2pt}{10pt}} and the domain
            \inblank{\blank{2.6cm}}{\gradedblank[id=q19h, type=short, points=0.5, answer={(-inf,inf)}, check=text, display={$(-\infty,\infty)$}]{2.6cm-2pt}{10pt}}. In one sentence, say why that domain is not like $g$'s.
 \rubricbox{\gradedblank[id=q19i, type=rubric, points=1, rubric={../rubrics/alg2-u05-actual-test-q19i.yaml}, display={index $3$ is odd: cubing keeps the sign, so no input is forbidden}]{\linewidth}{10pt}}

  \item \textbf{Solve. Test every candidate in the original equation.}
        \\[6pt]
        \answerbox[Answer:]{\gradedblank[id=q20a, type=short, points=2, answer={4}, check=numeric, display={$x=4$ only}]{1.1in}{10pt}}(a) $\sqrt{3x+4}=x$
        \\
        \ans{Square: $3x+4=x^2$, so $x^2-3x-4=0$ and $(x-4)(x+1)=0$, giving candidates $x=4$ and
        $x=-1$. Check $x=4$: $\sqrt{16}=4$ \checkmark. Check $x=-1$: $\sqrt1=1$, but the right side is
        $-1$, so $x=-1$ is \textbf{extraneous}. Solution: $x=4$ only. (Note the tell --- the
        non-radical side is negative at the rejected value.)}
        \\
        \answerbox[Answer:]{\gradedblank[id=q20b, type=short, points=1, answer={-7}, check=numeric, display={$x=-7$}]{1.1in}{10pt}}(b) $\sqrt[3]{x-1}=-2$
        \\
        \ans{Cube both sides --- legal in both directions, since every real number has exactly one
        real cube root: $x-1=-8$, so $x=-7$. Check: $\sqrt[3]{-8}=-2$ \checkmark. No extraneous risk
        with an odd index.}
        \\
        \answerbox[Answer:]{\gradedblank[id=q20c, type=short, points=1, answer={no solution}, check=text, display={no real solution}]{1.1in}{10pt}}(c) $\sqrt{2x-5}+7=3$
        \\
        \ans{Isolate first: $\sqrt{2x-5}=-4$. An isolated \emph{even} root is never negative, so
        \textbf{no real solution} --- and no algebra was needed. (Squaring anyway gives
        $x=\frac{21}{2}$, a phantom that fails the check.)}
        \\
        \answerbox[Answer:]{\gradedblank[id=q20d, type=short, points=1, answer={64, -64}, check=set, display={$x=\pm64$}]{1.1in}{10pt}}(d) $x^{2/3}=16$
        \\
        \ans{Raise both sides to the reciprocal power $\frac32$: $x=\pm 16^{3/2}=\pm64$. The
        \emph{numerator} $2$ is even --- it is a square in disguise --- so write the $\pm$ yourself.
        Both check: $(\pm64)^{2/3}=(\pm4)^2=16$ \checkmark}

  \item \textbf{Composition.} Let $f(x)=\sqrt{x}$ and $g(x)=x-9$.
        \\[4pt]
        (a) $(f\circ g)(x)=$ \inblank{\blank{3.0cm}}{\gradedblank[id=q21a, type=short, points=0.5, answer={sqrt(x-9)}, display={$\sqrt{x-9}$}]{3.0cm-2pt}{10pt}}, \quad domain: \inblank{\blank{2.6cm}}{\gradedblank[id=q21b, type=short, points=0.5, answer={x>=9}, display={$x\ge9$}]{2.6cm-2pt}{10pt}}
        \\[4pt]
        (b) $(g\circ f)(x)=$ \inblank{\blank{3.0cm}}{\gradedblank[id=q21c, type=short, points=0.5, answer={sqrt(x)-9}, display={$\sqrt{x}-9$}]{3.0cm-2pt}{10pt}}, \quad domain: \inblank{\blank{2.6cm}}{\gradedblank[id=q21d, type=short, points=0.5, answer={x>=0}, display={$x\ge0$}]{2.6cm-2pt}{10pt}}
        \\*[4pt]
        (c) Are these the same function? Answer using \emph{both} the rule and the domain.
 \rubricbox{\gradedblank[id=q21e, type=rubric, points=1, rubric={../rubrics/alg2-u05-actual-test-q21e.yaml}, display={No: different rules (right $9$ vs down $9$) and different domains ($x\ge9$ vs $x\ge0$)}]{\linewidth}{10pt}}\par
        (d) $(f\circ g)(25)=$ \inblank{\blank{2.2cm}}{\gradedblank[id=q21f, type=short, points=0.5, answer={4}, check=numeric, display={$\sqrt{16}=4$}]{2.2cm-2pt}{10pt}}
        \\[4pt]
        (e) For $p(x)=3x-2$ and $q(x)=x^2+1$, find $(p\circ q)(-3)$. \inblank{\blank{3.0cm}}{\gradedblank[id=q21g, type=short, points=0.5, answer={28}, check=numeric, display={$q(-3)=10$, then $p(10)=28$}]{3.0cm-2pt}{10pt}}

  \item \textbf{Inverses.}
        \\[4pt]
        (a) $f(x)=5x+3$. Find $f^{-1}(x)$. \inblank{\blank{4.0cm}}{\gradedblank[id=q22a, type=short, points=1, answer={(x-3)/5}, display={$f^{-1}(x)=(x-3)/5$}]{4.0cm-2pt}{10pt}}
        \\*[4pt]
        (b) Verify by composition --- \emph{both} orders --- that your answer really is the inverse. \rubricbox{\gradedblank[id=q22b, type=rubric, points=2, rubric={../rubrics/alg2-u05-actual-test-q22b.yaml}, display={$f(f^{-1}(x))=x$ and $f^{-1}(f(x))=x$: both orders return $x$}]{\linewidth}{34pt}}\par
        \answerbox[$g^{-1}(x)=$]{\gradedblank[id=q22c, type=short, points=1, answer={x^2+4}, display={$x^2+4$}]{0.85in}{10pt}\quad\textbf{domain:}\ \gradedblank[id=q22d, type=short, points=0.5, answer={x>=0}, display={$x\ge0$}]{0.65in}{10pt}}(c) $g(x)=\sqrt{x-4}$. Find $g^{-1}(x)$ and state its domain.
        \\
        \ans{Swap and solve: $x=\sqrt{y-4}$, so $x^2=y-4$ and $y=x^2+4$. Domain: $x\ge0$, because
        that is the \emph{range} of $g$ --- read it off, never guess it.}
        \\
        (d) The point $(13,3)$ lies on the graph of $g$. What point must lie on the graph of
            $g^{-1}$? \inblank{\blank{2.6cm}}{\gradedblank[id=q22e, type=short, points=0.5, answer={(3,13)}, check=text]{2.6cm-2pt}{10pt}}
\end{enumerate}

\begin{teachernote}
Item 1(f) is the Unit 5 habit-breaker: ``the graph stops'' is a \emph{complete} answer to the
left-end question, not a missing one --- do not accept ``$-\infty$'' or a blank. Item 1(g): the
endpoint is included, so $[-4,-3)$; a student writing $(-4,-3)$ has forgotten that $\sqrt0$ is real.
Item 2(e): $x^{5/4}$ alone earns most of the credit, but the item asked for a radical. Item 3: an
answer that is simplified but not \emph{fully} simplified is half an answer; 3(b) catches the
students who look for perfect squares under a cube root. Item 4(d): watch for the $a^2-b^2$ sign
flip in the denominator. Item 5(d) is the most-missed --- a negative $a$ makes the anchor a
\textbf{maximum}, and the range flips with it; 5(f) must name the \emph{index}, not just assert
``all reals.'' Item 6(a) is the graded check: ``$x=4$ and $x=-1$'' as a final answer earns no credit,
because the check is the step being assessed --- and 6(c) should be answered \emph{before} any
squaring. Item 7(c) must cite the domain, not only the rule. Item 8(b) is A2.F.2i's justification
clause --- both orders, or it is not a verification.
\end{teachernote}

% ======================================================================
\boxguard[36]% the heading, its instructions and the first item start together
\parthead{Part D --- Extended Response (12 pts)}
Explain your reasoning in full sentences.

\begin{enumerate}[leftmargin=1.9em, itemsep=16pt]
  \boxguard[30]
  \item \textbf{The full analysis.} Consider $f(x)=\sqrt{x-4}+1$.
        \textbf{(a)} Give the anchor point, the domain, and the range, both in interval notation.
        \textbf{(b)} Find the $x$-intercept and the $y$-intercept, or explain why one does not exist.
        \textbf{(c)} List the transformations that take $y=\sqrt{x}$ to $f$.
        \textbf{(d)} Find $f^{-1}(x)$ and state its domain, saying exactly where that restriction
        comes from.
        \textbf{(e)} Justify with composition --- both orders --- that $f$ and $f^{-1}$ are inverses.
        \textbf{(f)} A classmate writes $y=(x-1)^2+4$ with \emph{no} restriction and calls it the
        inverse. Explain why that is not the inverse of $f$.
        \rubricbox{\gradedblank[id=q23, type=rubric, points=6, rubric={../rubrics/alg2-u05-actual-test-q23.yaml}, display={(a) Anchor $(4,1)$; domain $[4,\infty)$; range $[1,\infty)$. Both start at the anchor.\\[3pt]
        (b) \textbf{Neither exists.} An $x$-intercept would need $\sqrt{x-4}+1=0$, i.e.
        $\sqrt{x-4}=-1$, and an even root is never negative. A $y$-intercept would need $f(0)$, but
        $0$ is not in the domain --- the graph does not reach the $y$-axis.\\[3pt]
        (c) Right $4$, up $1$. ($a=1$: no stretch, no reflection.)\\[3pt]
        (d) Swap and solve: $x=\sqrt{y-4}+1$, so $x-1=\sqrt{y-4}$, $(x-1)^2=y-4$, and
        $f^{-1}(x)=(x-1)^2+4$ for $x\ge1$. The restriction is the \textbf{range of $f$}, written
        down: whatever $f$ can output is exactly what $f^{-1}$ can accept.\\[3pt]
        (e) $f\!\left(f^{-1}(x)\right)=\sqrt{(x-1)^2+4-4}+1=\sqrt{(x-1)^2}+1=|x-1|+1$, and since
        $x\ge1$ that is $(x-1)+1=x$. Going the other way,
        $f^{-1}\!\left(f(x)\right)=\left(\sqrt{x-4}+1-1\right)^2+4=\left(\sqrt{x-4}\right)^2+4=(x-4)+4=x$
        for $x\ge4$. Both orders return $x$, so the two functions undo each other.\\[3pt]
        (f) Without the restriction it is the \emph{whole} parabola --- two arms --- and it fails the
        horizontal line test, so it cannot be the inverse of anything: it sends both $x=0$ and $x=2$
        to $5$. It also accepts inputs $f$ never produces; $f$ never outputs a number below $1$, so
        $x=0$ has no business being fed to $f^{-1}$. Only the right arm, $x\ge1$, is the inverse.}]{\linewidth}{6.2cm-6pt}}

  \boxguard[30]

  \item \textbf{Modeling.} Highway engineers estimate the greatest safe speed on a flat curved ramp
        of radius $r$ feet as
        \[ s(r)=4\sqrt{r}\ \text{ miles per hour.} \]
        \textbf{(a)} State the domain that makes sense here, and connect it to the rule about what an
        even index will accept.
        \textbf{(b)} Find $s(25)$ and $s(100)$, with units. What does quadrupling the radius do to the
        safe speed, and which feature of $\sqrt{\,\cdot\,}$ is responsible?
        \textbf{(c)} A ramp is posted at $36$ miles per hour. Find its radius algebraically, and check
        your answer.
        \textbf{(d)} Write the inverse, $r$ as a function of $s$, and say what it tells an engineer
        that $s(r)$ does not.
        \textbf{(e)} Part (c) required squaring both sides. Explain why there is no danger of an
        extraneous solution in this problem.
        \rubricbox{\gradedblank[id=q24, type=rubric, points=6, rubric={../rubrics/alg2-u05-actual-test-q24.yaml}, display={(a) $r\ge0$, i.e. $[0,\infty)$. The index is even, so the radicand may not be negative ---
        and the context says the same thing twice over, since a curve of negative radius is not a
        thing. Here the algebra and the situation agree; that is not always true.\\[3pt]
        (b) $s(25)=4\sqrt{25}=4\cdot5=20$ miles per hour; $s(100)=4\sqrt{100}=4\cdot10=40$ miles per
        hour. Quadrupling the radius only \textbf{doubles} the safe speed, because
        $\sqrt{4r}=2\sqrt{r}$ --- an exponent of $\frac12$ flattens growth, which is the shape of the
        square root graph.\\[3pt]
        (c) $4\sqrt{r}=36$, so $\sqrt{r}=9$ and $r=81$ feet. Check: $s(81)=4\sqrt{81}=4\cdot9=36$
        \checkmark\\[3pt]
        (d) Solve for $r$: $\sqrt{r}=\frac{s}{4}$, so $r=\dfrac{s^2}{16}$ for $s\ge0$. The original
        answers ``how fast is this ramp safe for?''; the inverse answers the question an engineer
        actually asks --- ``how wide must I build the curve to post a given speed?''\\[3pt]
        (e) An extraneous root only appears when squaring hides a \emph{sign} disagreement --- the
        step $a^2=b^2$ forgets whether $a=b$ or $a=-b$. Here both sides are non-negative:
        $\sqrt{r}\ge0$ and a posted speed of $36$ is positive, so there is no hidden sign to lose and
        the single candidate must be genuine. Checking it is still good practice.}]{\linewidth}{5.5cm-6pt}}
\end{enumerate}

\begin{teachernote}
\textbf{Scoring Part D (6 pts each).} D1: 1 pt (a) anchor with \emph{both} interval-notation answers;
1 pt (b) \emph{both} non-existences with reasons --- ``no $x$-intercept'' with no reason earns half;
1 pt (c) both transformations; 1 pt (d) the rule \emph{and} $x\ge1$ \emph{and} the sentence that it is
the range of $f$ (no credit for the restriction if it is guessed as $x\ge0$, the reflex answer);
1 pt (e) both compositions --- one order is not a verification, and this is the standard's own
``justify'' clause; 1 pt (f) the horizontal line test or the two-arms argument. A student who writes
$(x-1)^2+4$ with no restriction in (d) and then cannot explain (f) has missed the unit's closing
idea --- score both accordingly. D2: 1 pt (a) $r\ge0$ tied to the even index; 1 pt (b) both values
with units \emph{and} ``doubles'' with $\sqrt{4r}=2\sqrt{r}$ (half credit for ``it doubles'' with no
reason); 1 pt (c) $r=81$ ft; 1 pt the check; 1 pt (d) $r=\frac{s^2}{16}$; 1 pt (e) the sign argument.
Do not accept ``because I checked it'' for (e) --- the question asks why the danger was never there.
\end{teachernote}

\end{document}
